\documentclass[11pt,a4paper]{article}

% ---------- packages ----------
\usepackage[utf8]{inputenc}
\usepackage[T1]{fontenc}
\usepackage{amsmath,amssymb,amsthm}
\usepackage{geometry}
\geometry{margin=2.4cm}
\usepackage{booktabs}
\usepackage{array}
\usepackage{enumitem}
\usepackage{xcolor}
\usepackage{siunitx}
\usepackage{graphicx}
\usepackage[colorlinks=true,linkcolor=black,citecolor=blue,urlcolor=blue]{hyperref}

% ---------- operators / macros ----------
\DeclareMathOperator{\sat}{sat}
\newcommand{\COP}{\mathrm{COP}}
\newcommand{\Qmax}{\dot{Q}_{\max}}
\newcommand{\Ttop}{T_{\mathrm{top}}}
\newcommand{\Tbot}{T_{\mathrm{bot}}}
\newcommand{\Ths}{T_{hs}}          % HP supply  (condenser outlet)
\newcommand{\Thr}{T_{hr}}          % HP return  (condenser inlet)
\newcommand{\Tsp}{T_{\mathrm{sp}}} % HP leaving-water setpoint (the control)
\newcommand{\Tair}{T_{\mathrm{air}}}
\newcommand{\Tamb}{T_{\mathrm{amb}}}
\newcommand{\Tretd}{T_{\mathrm{ret},d}}
\newcommand{\Tsd}{T_{s,d}}         % delivered district supply temp
\newcommand{\Tsdref}{T_{s,d}^{\star}}
\newcommand{\mh}{\dot m_h}         % HP condenser-loop flow
\newcommand{\md}{\dot m_d}         % district draw flow
\newcommand{\Qhp}{Q_{\mathrm{hp}}}
\newcommand{\Pel}{P_{\mathrm{el}}}
\newcommand{\Vtop}{V_{\mathrm{top}}}
\newcommand{\cw}{c}                % specific heat of water
\newcommand{\code}[1]{\texttt{#1}}

\title{\textbf{Grey-box Model of the Profondeville Energy Center\\for Demand-Response MPC}}
\author{Karno / Profondeville (K-0001)}
\date{\today}

\begin{document}
\maketitle

\begin{abstract}
This note derives a control-oriented grey-box model of the Profondeville district
heating energy center, for the purpose of a demand-response (DR) model predictive
controller (MPC) driven by electricity price, PV production and a district
consumption forecast. The physical boundary of the model is the \emph{supply
temperature to the district}; the apartments are not modelled. The controllable
subsystem consists of the air-source heat pump (ASHP, modulating), the geothermal
heat pump (GSHP, on/off) and the distribution 3-way mixing valve. We present (i) a
detailed plant model suitable for system identification, keeping the two internal
PID loops explicit, and (ii) a reduced model suitable for online optimisation, in
which the PID loops collapse into algebraic set-point tracking plus feasibility
constraints. A single continuous input (the ASHP) and a single storage state (the
buffer tank) are sufficient for a first controller; the geothermal on/off unit and
a stratification (thermocline) state are documented as upgrades.
\end{abstract}

\tableofcontents

% ==================================================================
\section{System and modelling scope}
% ==================================================================

In winter the energy center produces hot water into a \SI{2.5}{m^3} buffer tank,
charged by two heat pumps, and feeds a district network whose supply temperature is
regulated by a 3-way mixing valve. The relevant energy flows are:

\begin{itemize}[nosep]
  \item \textbf{ASHP} (air source, $\sim$\SI{20}{kW}): inverter-\emph{modulating},
        $0$--$100\,\%$. Its onboard controller regulates the leaving-water
        temperature to a set-point. This is the primary continuous actuator.
  \item \textbf{GSHP} (ground source, $\sim$\SI{40}{kW}): \emph{on/off} only.
        Treated in v1 either as off or as a measured exogenous heat input.
  \item \textbf{Buffer tank} (\SI{2.5}{m^3}): the thermal storage that provides
        the demand-response flexibility.
  \item \textbf{3-way valve} (\code{FCV\_702}): blends hot tank water with the
        warm district return to hold the district supply temperature on a
        weather-compensated heating curve.
\end{itemize}

\paragraph{Measurement reality (design constraint).}
Water \emph{flows} and \emph{heat meters} on the heat-pump loops are \emph{not}
currently instrumented and are assumed unavailable. Only temperatures, valve/pump
positions, the district flow \code{FQT\_701\_OUT}, and heat-pump \emph{electrical}
power are measured. The model is therefore built so that all unmeasured flows are
absorbed into identifiable parameters, and the heat-pump heat input is inferred
from measured electrical power.

\paragraph{Control architecture (cascade).}
The MPC does not command power directly. It commands a \emph{set-point}; the
electrical power and the delivered heat are \emph{responses} of the heat pump:
\begin{equation}
  \underbrace{u=\Tsp\ (+\,\text{enable }\delta)}_{\text{MPC decision}}
  \;\xrightarrow{\ \text{HP map}\ }\;
  (\Pel,\ \Qhp)
  \;\xrightarrow{\ \text{tank}\ }\;
  (\Ttop,\Tbot)
  \;\xrightarrow{\ \text{valve}\ }\;
  \Tsd .
\end{equation}

% ==================================================================
\section{Nomenclature and measured signals}
% ==================================================================

\renewcommand{\arraystretch}{1.15}
\begin{table}[h]
\centering
\small
\begin{tabular}{@{}l l l l@{}}
\toprule
Symbol & Meaning & Unit & BigQuery column \\
\midrule
\multicolumn{4}{@{}l}{\emph{States}}\\
$\Ttop$ & Tank top temperature & \si{\celsius} & \code{TT\_601} \\
$\Tbot$ & Tank bottom temperature & \si{\celsius} & \code{TT\_602} \\
$\Ths$  & ASHP supply (condenser outlet) & \si{\celsius} & \code{PAC\_501\_T\_OUT\_EAU\_ECHANGEUR} \\
$\Thr$  & ASHP return (condenser inlet)  & \si{\celsius} & \code{PAC\_501\_T\_IN\_EAU\_ECHANGEUR} \\
$\Vtop$ & Hot-zone volume (optional, thermocline) & \si{m^3} & --- (observer) \\
\midrule
\multicolumn{4}{@{}l}{\emph{Control inputs}}\\
$\Tsp$  & ASHP leaving-water set-point & \si{\celsius} & \code{PAC\_501\_CONSIGNE\_ACTUEL} \\
$\delta$ & ASHP enable (on/off) & --- & \code{PAC\_501\_RUN}/\code{\_LIBERATION} \\
$\theta$ & 3-way valve position & --- & \code{FCV\_702\_POSITION} \\
$\delta_g$ & GSHP enable (v1.5) & --- & \code{PAC\_301\_RUN} \\
\midrule
\multicolumn{4}{@{}l}{\emph{Measured disturbances / forecasts}}\\
$\Tair$ & Outdoor air temperature & \si{\celsius} & \code{TT\_Text} \\
$Q_{\mathrm{dist}}$ & District heat load (forecast; tank driver) & \si{kW} & forecast / $\md\cw\Delta T$ \\
$\md$   & District draw flow (historical only) & \si{m^3/h} & \code{FQT\_701\_OUT} \\
$\Tretd$& District return temperature & \si{\celsius} & \code{TT\_702} \\
$\Tamb$ & Plant-room ambient temperature & \si{\celsius} & (assumed / const) \\
$\lambda(t)$ & Electricity price & \si{EUR/kWh} & external \\
$\mathrm{PV}(t)$ & PV production / forecast & \si{kW} & external \\
\midrule
\multicolumn{4}{@{}l}{\emph{Measured outputs}}\\
$\Pel$  & ASHP electrical power & \si{kW} & \code{PAC\_501\_PUISSANCE} / \code{EM\_PAC\_501\_PWR\_TOT\_P} \\
$\Tsd$  & District supply temperature & \si{\celsius} & \code{TT\_701} \\
\bottomrule
\end{tabular}
\caption{Model variables and the logged signals they map to. All parameters
($C_\bullet, \kappa, UA_\bullet, \alpha, \beta,\dots$) are identified; unmeasured
flows are folded into them.}
\end{table}

Constants: $\cw$ = specific heat of water (\si{kJ/(kg.K)}), $\rho$ = density.
Flows are expressed as mass flows $\dot m = \rho\,\dot V$.

% ==================================================================
\section{Heat-pump model (ASHP)}
% ==================================================================

\subsection{Internal set-point-tracking loop (PID $\to$ modulation)}
The inverter modulates the compressor, $m\in[0,1]$, to drive the leaving-water
temperature $\Ths$ to its set-point $\Tsp$. This is the first of the two internal
PID loops:
\begin{align}
  e(t) &= \Tsp(t) - \Ths(t), \\
  m(t) &= \sat_{[0,1]}\!\Big( K_p\,e(t) + K_i\!\int_0^t e(\tau)\,\mathrm d\tau \Big).
  \label{eq:hp-pid}
\end{align}

\begin{quote}\small\itshape
The gains $K_p,K_i$ are \textbf{not identified}. Equation~\eqref{eq:hp-pid} is
documentation of the onboard controller, not a block that is fitted. The data
already contains this loop's \emph{output}: the modulation, $\Pel$ and $\Ths$ are
all logged. For identification (\S\ref{sec:sysid}) we regress on those measured
signals; for control (\S\ref{sec:mpc}) we collapse the loop to $\Ths=\Tsp$. The
gains enter neither.
\end{quote}
The modulation delivers thermal power to the condenser water, capped by the
air-temperature-dependent capacity ceiling:
\begin{equation}
  \Qhp(t) = m(t)\,\Qmax\!\big(\Tair(t)\big), \qquad
  \Qmax(\Tair) \approx a_0 + a_1\,\Tair .
  \label{eq:qhp}
\end{equation}

\subsection{Electrical power and COP}
The electrical power (the quantity multiplied by price) follows from the coefficient
of performance, which decreases with lift (sink temperature $\Ths$ up, source
temperature $\Tair$ down):
\begin{equation}
  \Pel(t) = \frac{\Qhp(t)}{\COP\big(\Tair,\Ths\big)}.
  \label{eq:pel}
\end{equation}
Two identifiable COP forms:
\begin{equation}
  \text{(Carnot-scaled)}\quad
  \COP = \eta\,\frac{\Ths+273.15}{\Ths-\Tair+\Delta T_{\text{pp}}},
  \qquad
  \text{(affine)}\quad
  \COP = b_0 + b_1\Tair + b_2\Ths .
  \label{eq:cop}
\end{equation}
Since $\Pel,\Tair,\Ths$ are all measured, \eqref{eq:pel}--\eqref{eq:cop} are fit by
regression on historical data. This is the only place nonlinearity enters the cost;
Section~\ref{sec:mpc} discusses convex handling.

\subsection{Condenser-loop dynamics (optional refinement)}
The heat-pump water loop draws from the tank bottom, is heated in the condenser and
returns to the tank top. The condenser circulator is \emph{fixed-speed}, so the loop
flow $\mh$ is a \textbf{single unknown constant}, not a signal. With loop thermal
capacitance $C_h$:
\begin{align}
  C_h\,\dot{\Ths} &= \Qhp - \mh\,\cw\,(\Ths-\Thr), \label{eq:ths}\\
  \tau_r\,\dot{\Thr} &= \Tbot - \Thr. \label{eq:thr}
\end{align}
Equation~\eqref{eq:thr} is a first-order lag capturing pipe/mixing transport from
the tank bottom to the condenser inlet; in the fast limit $\Thr\to\Tbot$.

\begin{quote}\small\itshape
On the unknown flow $\mh$. It is never needed in physical units. (i) In
\eqref{eq:ths} it only appears as the ratio $\mh\cw/C_h$, one lumped coefficient
identified from the $\Ths$ transient. (ii) For the \emph{heat delivered to the
tank} we do not use \eqref{eq:ths} at all: we take it from the measured electrical
power, $\Qhp=\alpha\,\Pel$ (\S\ref{sec:tank}). Consequently the tank model below
contains no $\mh$: the flow's only remaining footprint is the recirculation
\emph{stirring} it causes, which lumps into a single mixing coefficient $\kappa$.
This condenser-loop block is therefore an \emph{optional} refinement; the reduced
model drops it entirely.
\end{quote}

% ==================================================================
\section{Geothermal heat pump (GSHP)}
\label{sec:gshp}
% ==================================================================

The GSHP ($\sim$\SI{40}{kW}) is \textbf{on/off only} and, in the present operation,
its start/stop is decided by the existing plant logic --- it is \emph{not} an MPC
decision variable in v1. It must nevertheless be modelled, for one reason: to fit
the tank on data segments where the GSHP was running, its heat input has to be
accounted for, otherwise the identifier attributes the GSHP's charging to the ASHP
or to the tank parameters and the fit is corrupted.

Because the ground-source temperature is nearly constant across a heating season,
the on-state heat output is essentially constant. The model is therefore a switched
constant source, or --- preferred, since it is faithful and needs no COP assumption
--- the \emph{measured} electrical power scaled by an identified proxy:
\begin{equation}
  Q_{\mathrm{gshp}}(t) =
  \begin{cases}
    \alpha_g\,P_{\mathrm{el},g}(t) & \text{(measured-power form, preferred)},\\[2pt]
    \delta_g(t)\,\bar Q_{\mathrm{gshp}} & \text{(switched-constant form)},
  \end{cases}
  \label{eq:gshp}
\end{equation}
with $\delta_g\in\{0,1\}$ the on/off state (\code{PAC\_301\_RUN}),
$P_{\mathrm{el},g}$ its electrical power (\code{EM\_PAC\_301\_PWR\_TOT\_P} /
\code{PAC\_301\_PUISSANCE}), and $\alpha_g$ or $\bar Q_{\mathrm{gshp}}$ identified.
This term enters the tank top node in \eqref{eq:ttop} exactly like the ASHP heat.
Structurally the two heat pumps are identical inputs to the tank; the difference is
only that $\Tsp$ (ASHP) is a decision variable whereas $\delta_g$ (GSHP) is, for
now, an exogenous \emph{measured disturbance}. If the GSHP is later added to the
controller, \eqref{eq:gshp} is already the required model (a binary control with a
constant heat quantum) --- no re-identification needed.

% ==================================================================
\section{Buffer-tank model}
\label{sec:tank}
% ==================================================================

\subsection{Heat inputs (measured, no HP flow needed)}
Both heat pumps charge the \emph{same} tank. Their heat inputs are taken from the
measured electrical powers, so the tank model contains no heat-pump flow:
\begin{equation}
  \Qhp = \alpha\,\Pel \quad(\text{ASHP}), \qquad
  Q_{\mathrm{gshp}} = \alpha_g\,P_{\mathrm{el},g}\quad(\text{GSHP, \S\ref{sec:gshp}}),
  \label{eq:heatin}
\end{equation}
with efficiency-proxy constants $\alpha,\alpha_g$ (identified). Both inject into the
top node.

\subsection{Two-node stratified model (v1)}
The tank is discretised into a top node ($\Ttop$) and a bottom node ($\Tbot$) of
fixed volumes. The district acts as a \emph{heat sink}: hot water is withdrawn from
the top and returns cold to the bottom, removing heat at the rate
$Q_{\mathrm{dist}}$. We drive the tank with this \emph{forecasted heat} directly
rather than with the draw flow $\md$, because the demand is forecast as a load (kW),
not a flow, and reconstructing $\md\,\cw\,\Delta T$ only to collapse it back to a heat
is needless --- the useful product is the heat. Since extraction is from the top, the
load lands on the top node, and with the district return sitting near the bottom node
($\Tretd\approx\Tbot$) the whole draw
$\md\,\cw(\Ttop-\Tbot)\approx\md\,\cw(\Ttop-\Tretd)=Q_{\mathrm{dist}}$ collapses onto
it. A single factor $\beta$ carries any constant mismatch between the forecast
boundary and this tank's outlet (network/standing losses, metering point, a shared
manifold); $\beta\approx1$ when $Q_{\mathrm{dist}}$ is the tank-outlet heat meter. The
heat pumps enter as the heat terms \eqref{eq:heatin}; their recirculation, conduction,
and the flow-driven inter-node transport appear as one mixing coefficient $\kappa$.
The node energy balances are:
\begin{align}
  C_{\mathrm{top}}\,\dot{\Ttop} =\;&
    \Qhp + Q_{\mathrm{gshp}}
    - \beta\,Q_{\mathrm{dist}}
    - \kappa\,(\Ttop-\Tbot)
    - UA_t\,(\Ttop-\Tamb), \label{eq:ttop}\\[2pt]
  C_{\mathrm{bot}}\,\dot{\Tbot} =\;&
    \kappa\,(\Ttop-\Tbot)
    - UA_b\,(\Tbot-\Tamb). \label{eq:tbot}
\end{align}
The whole tank then loses exactly $\beta\,Q_{\mathrm{dist}}$ to the district (plus the
standing losses), which is energy-consistent by construction. Here $\kappa$ (W/K) is
the mixing/conduction coefficient --- it absorbs the fixed-speed HP recirculation and
the flow-driven transport, and produces the gradual $\Ttop$ droop of a non-ideal tank
--- and $UA_t,UA_b$ are standing-loss coefficients. \textbf{Identified parameters:}
$C_{\mathrm{top}},C_{\mathrm{bot}},\kappa,UA_t,UA_b,\alpha,\alpha_g,\beta$; the tank
capacitances are anchored near their geometric values $\rho V_\bullet\cw$ (the buffer
volume is known), which keeps $\beta$ from trading off against $C_{\mathrm{top}}$.
\emph{No heat-pump flow $\mh$ appears.}

\emph{Bottom-node anchor.} Dropping the explicit cold-return term removes the bottom
node's cold source, so $\Tbot$ can drift slightly warm via $\kappa$. If the bottom
fit matters, retain the measured return as a weak pull $-\,\gamma\,(\Tbot-\Tretd)$ in
\eqref{eq:tbot} ($\Tretd=\code{TT\_702}$, a measured input); for the top-node dynamics
that drive supply and DR it is immaterial.

\paragraph{Opposite-curvature diagnostic (when is two-node enough?).}
The fixed volumes $C_{\mathrm{top}},C_{\mathrm{bot}}$ are the modelling
approximation: physically the stored energy lives in a \emph{hot volume} and a
\emph{cold volume} whose boundary (the thermocline) moves, so charge/discharge is
really a change of \emph{volume} at nearly constant temperatures. Freezing the
volumes forces that volume motion to be re-expressed as a \emph{temperature} change
of fixed nodes; the coefficient $\kappa$ is the artifact that absorbs it. This
substitution has a characteristic signature in the discharge shape of $\Ttop$:
\begin{itemize}[nosep]
  \item \textbf{Two-node model:} $\Ttop$ relaxes \emph{exponentially} toward
        $\Tbot$ --- fastest drop first, then tapering (convex, decreasing slope).
  \item \textbf{Real moving thermocline:} $\Ttop$ holds a \emph{plateau} while the
        front is below the sensor, then drops sharply when it arrives (concave,
        flat-then-steep).
\end{itemize}
These are \emph{opposite curvatures}. Hence an observed ``flat for a while, then
lower'' discharge is the thermocline signature and the regime the two-node model is
structurally worst at. Use this as a quantitative trigger: fit the two-node model
first and inspect its residuals around the deepest discharges. If the model
consistently \emph{leads} the data --- $\Ttop$ predicted to droop while the
measurement is still on its plateau, then catching up after the real cliff --- the
knee is too sharp for two nodes and the $\Vtop$ thermocline state
(\S\ref{sec:vtop}) is required. For a DR controller this residual pattern is not
cosmetic: it means the model under-predicts available flexibility on the plateau
and then over-predicts it past the cliff, i.e.\ it mistimes exactly the moment the
tank runs out. A gentle knee (soft plateau, mild drop) stays inside the two-node
regime and $\kappa$ can absorb it.

\paragraph{Validity.}
A two-node model reproduces a \emph{smooth} $\Ttop$ decay during discharge
\emph{by construction}: it cannot represent a sharp thermocline sweeping past the
sensor. It is therefore valid only in the regime where the real tank actually
de-stratifies smoothly. This should be checked on the deepest observed discharges;
where $\Ttop$ instead drops abruptly (a ``cliff''), use the thermocline state below.

\subsection{Thermocline state (v1.5, optional)}
\label{sec:vtop}
The point sensors $\Ttop,\Tbot$ sit near the tank ends. Their value is not the
usable stored energy: a thin hot layer over a large cold volume stores little.
Introduce the \emph{hot-zone volume} $\Vtop$ (position of the thermocline). Because
the HP loop is not flow-metered, its \emph{volumetric} charging rate is recovered
from the measured \emph{electrical} power via the heat balance:
\begin{equation}
  \dot V_{\text{charge}}
  = \frac{\Qhp}{\rho\,\cw\,(\Ttop-\Tbot)}
  = \frac{\alpha\,\Pel}{\rho\,\cw\,(\Ttop-\Tbot)},
  \qquad \Qhp \approx \alpha\,\Pel,
  \label{eq:vcharge}
\end{equation}
while the discharge draws hot water from the top at the measured district flow times
the fraction pulled from the tank by the valve, $g(\theta)$:
\begin{equation}
  \boxed{\;
  \dot{\Vtop} = \frac{\alpha\,\Pel}{\rho\,\cw\,(\Ttop-\Tbot)}
              - \md\,g(\theta)\; }
  \label{eq:vtop}
\end{equation}
Equation~\eqref{eq:vtop} degenerates as $\Ttop-\Tbot\to0$ (no distinct thermocline);
clamp/freeze $\Vtop$ below a $\Delta T$ threshold.

\paragraph{Observability of $\Vtop$ from point sensors.}
$\Vtop$ is \emph{not} continuously measured, but it is partially observed as a
\emph{threshold event}: while the thermocline is below the top sensor, $\Ttop$
saturates at the hot value ($\Vtop \ge V_{\text{sensor}}$, exact value unknown);
the instant $\Ttop$ drops is a hard fix that the front has reached the sensor
height. This yields a predict/correct observer:
\begin{align}
  \text{predict:}\quad & \dot{\Vtop}\ \text{by }\eqref{eq:vtop}\ \text{(open-loop integration)},\\
  \text{correct:}\quad & \Vtop \leftarrow V_{\text{sensor}}\ \text{when a }\Ttop\ \text{(or }\Tbot\text{) crossing is detected.}
\end{align}
Between the two end-sensors the state is unobserved and relies on the integration;
the crossings cancel its drift.

% ==================================================================
\section{Distribution valve and heating curve}
\label{sec:valve}
% ==================================================================

The 3-way valve blends tank-top water with the district return to deliver the
supply temperature:
\begin{equation}
  \Tsd = \theta\,\Ttop + (1-\theta)\,\Tretd, \qquad \theta\in[0,1]. \label{eq:blend}
\end{equation}
A second internal PID drives $\theta$ so that $\Tsd\to\Tsdref$. \textbf{We do not
model this PID and do not identify its gains}: the valve is not a demand-response
actuator (it is absent from the MPC decision variables in \eqref{eq:obj}), so its
dynamics are irrelevant to the flexibility problem. The \emph{only} thing this loop
contributes to our model is its steady-state result --- the algebraic feasibility
constraint \eqref{eq:feasibility} --- because whenever the tank is hot enough the
valve holds $\Tsd=\Tsdref$ exactly, and when it is not the valve saturates. That is
all the MPC needs from the valve.

The set-point $\Tsdref$ is the weather-compensated heating curve (encoded in the
winter scenario, columns \code{MODE\_HIV\_ConsigneSupplyIfTempExt*Param}); a
two-level day/night form is
\begin{equation}
  \Tsdref(\Tair) =
  \begin{cases}
    T_{\text{cold}}, & \Tair < \bar T - \tfrac{\Delta}{2},\\
    T_{\text{mild}}, & \Tair > \bar T + \tfrac{\Delta}{2},
  \end{cases}
  \label{eq:heatcurve}
\end{equation}
with switch threshold $\bar T$ (\code{MODE\_HIV\_paramTemperatureExt}) and hysteresis
$\Delta$ (\code{paramOffsetTempExt}).

\paragraph{Feasibility (the key coupling).}
Blending \eqref{eq:blend} can only \emph{raise} the return toward $\Ttop$; the
valve saturates at $\theta=1$ once $\Ttop=\Tsdref$. Hence the deliverable-supply
constraint
\begin{equation}
  \Tsd = \Tsdref \iff \Ttop \ge \Tsdref(\Tair).
  \label{eq:feasibility}
\end{equation}
The right-hand inequality is the \emph{floor} of the usable storage: once $\Ttop$
falls to $\Tsdref$, flexibility is exhausted and the heat pumps must recharge.

% ==================================================================
\section{Reduced model and MPC formulation}
\label{sec:mpc}
% ==================================================================

\subsection{Collapsing the PID loops}
The low-level PIDs (HP loop \eqref{eq:hp-pid} and the valve loop of \S\ref{sec:valve}) are not re-optimised
by the MPC. On the MPC time step they are assumed to reach their set-points whenever
physically feasible, which replaces the fast dynamics by algebraic relations plus
saturation limits:
\begin{align}
  \Ths &= \Tsp & &\text{while } m<1 \ \ (\Qhp<\Qmax(\Tair)), \label{eq:collapse1}\\
  \Tsd &= \Tsdref(\Tair) & &\text{while } \Ttop\ge \Tsdref(\Tair). \label{eq:collapse2}
\end{align}
The heat-pump heat input, expressed directly from the decision variable (needed
because in a forward prediction $\Pel$ is an output, not yet known):
\begin{equation}
  \Qhp = \min\!\Big(\ G_h\,(\Tsp-\Tbot),\ \ \Qmax(\Tair)\ \Big),
  \qquad G_h \equiv \mh\,\cw,
  \label{eq:qhp-red}
\end{equation}
where $G_h$ (W/K) is the \emph{single lumped condenser conductance} of
\S\ref{sec:tank}: the unknown fixed pump flow times $\cw$, identified as one
constant. The two limbs are the compressor working ($\Ths=\Tsp$) and the
compressor saturated at full modulation.

\subsection{Reduced dynamics}
Dropping the fast condenser states \eqref{eq:ths}--\eqref{eq:thr}, the control model
keeps the two tank states (optionally $\Vtop$):
\begin{align}
  C_{\mathrm{top}}\,\dot{\Ttop} &=
     \Qhp + Q_{\mathrm{gshp}} - \beta\,Q_{\mathrm{dist}} - \kappa(\Ttop-\Tbot) - UA_t(\Ttop-\Tamb), \label{eq:red-top}\\
  C_{\mathrm{bot}}\,\dot{\Tbot} &=
     \kappa(\Ttop-\Tbot) - UA_b(\Tbot-\Tamb), \label{eq:red-bot}\\
  \Pel &= \Qhp\,/\,\COP(\Tair,\Tsp). \label{eq:red-pel}
\end{align}

\subsection{Constraints (where the PID logic becomes explicit)}
\begin{align}
  \Ttop &\ge \Tsdref(\Tair) && \text{valve saturation} \eqref{eq:feasibility}: \text{DR floor},\\
  \Ttop &\le T_{\max} && \text{tank ceiling},\\
  \Qhp &\le \Qmax(\Tair) && \text{compressor saturation } (m=1),\\
  \Tsp &\in [\Tsp^{\min},\Tsp^{\max}], \quad \delta\in\{0,1\} && \text{actuator limits}.
\end{align}

\subsection{Objective}
Minimise energy cost net of PV self-consumption over the horizon $H$:
\begin{equation}
  \min_{\{\Tsp(t),\,\delta(t)\}}\ \int_{0}^{H}
    \Big[\ \lambda(t)\,\Pel(t) \;-\; c_{\text{pv}}\,\min\!\big(\Pel(t),\mathrm{PV}(t)\big)\ \Big]\,\mathrm dt .
  \label{eq:obj}
\end{equation}

\subsection{Convexity note}
The plant is nonconvex (bilinear tank advection; $\Pel=\Qhp/\COP$). Two routes:
\begin{enumerate}[nosep]
  \item \textbf{Nonlinear MPC:} solve \eqref{eq:red-top}--\eqref{eq:obj} directly
        as an NLP (e.g.\ CasADi + IPOPT). Recommended for the first closed loop.
  \item \textbf{Convex approximation:} freeze the COP on the \emph{forecast}
        exogenous signals, $\COP(t)=\COP(\Tair(t),\Tsdref(t))$, so the cost
        $\lambda(t)/\COP(t)\cdot\Qhp$ becomes linear in $\Qhp$; with a fixed-node
        tank \eqref{eq:red-top}--\eqref{eq:red-bot} the problem is an LP/QP. This
        ignores the efficiency penalty of charging hotter, to be reintroduced by
        piecewise-linear COP later.
\end{enumerate}

% ==================================================================
\section{System identification with \texttt{ctsmTMB}}
\label{sec:sysid}
% ==================================================================

Identification is carried out with \texttt{ctsmTMB} (continuous--discrete
stochastic state-space models; maximum likelihood via an extended/unscented Kalman
filter). The model is a set of It\^o SDEs for the states plus discrete observation
equations for the sensors:
\begin{align}
  \mathrm d x &= f(x,u,t;\theta)\,\mathrm dt + \sigma(\theta)\,\mathrm d\omega,
  & &\text{(system, process noise $\omega$)} \label{eq:sde}\\
  y_k &= h(x_k,u_k;\theta) + e_k, \quad e_k\sim\mathcal N(0,\Sigma_e),
  & &\text{(observation, measurement noise)} \label{eq:obs}
\end{align}
and $\theta$ is estimated by maximising the Kalman-filter likelihood. The process
noise term is the reason to prefer this over plain least squares: it lets the
filter absorb the \emph{structural} error of the model (e.g.\ the thermocline
approximation of \S\ref{sec:tank}) instead of biasing the parameters into it.

\subsection{The reframing: the heat pumps enter as measured inputs, not maps}
The apparent obstacle is that the ``HP static map'' \eqref{eq:qhp}--\eqref{eq:cop}
is nonlinear, and $\Qhp$ is unmeasured. Both dissolve once one observes that
\textbf{during identification the electrical power is logged, not predicted}. We
never turn a set-point into power here; the machine already did that, and the
result is the measured column $\Pel$. Hence the heat pumps enter the tank SDE as
\emph{measured inputs} scaled by one constant each:
\begin{equation}
  \Qhp(t)=\alpha\,\underbrace{\Pel(t)}_{\text{input}}, \qquad
  Q_{\mathrm{gshp}}(t)=\alpha_g\,\underbrace{P_{\mathrm{el},g}(t)}_{\text{input}} .
  \label{eq:heat-inputs}
\end{equation}
There is no $\Qmax$, no $\sat$, no COP curve in the identification model: all of
that is already baked into the measured $\Pel$. The parameters $\alpha,\alpha_g$
\emph{are} the effective COPs, and the tank temperatures respond in proportion to
the injected heat, so they are identifiable. \emph{The internal PID gains
$K_p,K_i$ never appear} --- the loop's output ($\Pel,\Ths$) is measured.

\subsection{Closed-loop caveat: the ASHP enable signal is generated by \texorpdfstring{$\Ttop$}{Ttop}}
\label{sec:closedloop}

Section~\ref{sec:sysid}.1 treats $\Pel$ (and, by extension, $\delta$) as a
\emph{measured input} to be fed into the tank SDE without further comment. On
the K-0001 data this needs a caveat: the existing (pre-MPC) plant logic that
generated the historical $\delta,\Pel$ record is not free/exogenous, it is a
\textbf{feedback law on $\Ttop$ itself} --- a Schmitt-trigger (hysteresis)
relay,
\begin{equation}
  \delta(t) =
  \begin{cases}
    1 & \Ttop(t) \text{ crosses below } \SI{31}{\celsius}\ \text{(enable)},\\
    0 & \Ttop(t) \text{ crosses above } \SI{35}{\celsius}\ \text{(disable)},\\
    \delta(t^-) & \text{otherwise (inside the hysteresis band)}.
  \end{cases}
  \label{eq:enable-law}
\end{equation}

\paragraph{Empirical evidence (11-day window, 2026-05-13 to 2026-05-24, 5 min sampling).}
\begin{itemize}[nosep]
  \item Raw correlation $\mathrm{corr}(\Pel,\Ttop) = -0.57$ --- the wrong sign
        for a term meant to \emph{heat} $\Ttop$.
  \item Pre-whitened cross-correlation (AR(20) prewhitening, Box--Jenkins):
        peak $\Pel\!\to\!\Ttop$ correlation is at \textbf{lag 0}, $r=-0.39$
        (instantaneous, negative --- not the delayed positive response direct
        heating would produce). By contrast $\Pel\!\to\!\Tbot$ peaks at lag
        $+1$ sample ($+5$ min), $r=+0.30$: small positive delay, correct sign,
        consistent with $\Tbot$ being a passive receiver of whatever heat is
        actually added.
  \item $\delta$ (\code{PAC\_501\_RUN}) correlates $0.95$ with $\Pel$
        (expected: same event) but \emph{also} $-0.48$ with $\Ttop$ --- the
        enable command itself is triggered by low $\Ttop$, so it carries the
        same endogeneity as $\Pel$.
  \item $\Tsp$ (\code{PAC\_501\_CONSIGNE\_ACTUEL} / logged as
        \code{T\_set,HP,air}) is \textbf{constant} within this window --- a
        two-level seasonal flag over the full $\sim$2.7-year log, not
        something varied independently for identification purposes. No clean
        exogenous proxy is available in the historical record.
\end{itemize}

\begin{quote}\small\itshape
Why this matters for Stage~1. The Kalman-filter MLE in \eqref{eq:sde}--\eqref{eq:obs}
implicitly treats the input trajectory as independent of the process noise
driving the state. When $\delta$ (and hence $\Pel$) is actually generated by
\eqref{eq:enable-law} --- a function of $\Ttop$, the very state being
estimated --- that independence fails, and the estimator can no longer tell
``the plant responds this way to heat'' apart from ``the old controller reacts
this way to temperature.'' A preliminary two-node fit on this data (constant
$\bar Q_{\mathrm{gshp}}$, $\delta$ and $\Pel$ used as raw measured inputs
exactly as in \eqref{eq:heat-inputs}) showed the signature of this: $C_{\mathrm{top}}$,
$C_{\mathrm{bot}}$, $\kappa$ and $\alpha$ all pinning at whatever bound they
were given rather than converging, and the two channels' noise parameters
($\sigma_{y_1}$ vs.\ $\sigma_{y_2}$) splitting to opposite bounds --- one
channel's noise absorbing the mismatch instead of a real fit.

Direct closed-loop identification (fitting the plant on closed-loop data,
input treated as given) is not automatically wrong --- it is provably
consistent under prediction-error methods \emph{provided} the model structure
is correct and the noise is correctly specified (Ljung, \emph{System
Identification: Theory for the User}). The practical reading here: closed-loop
data is much less forgiving of structural misspecification (wrong injection
node, missing dynamics, oversimplified coupling) than open-loop data, because
feedback narrows the input's information content and concentrates any
structural error into biased parameter estimates rather than averaging it out.
\end{quote}

\paragraph{Fix: fold the known switching law into the SDE.}
Because \eqref{eq:enable-law} is fully specified (even though the in-band
modulation level is not), it should be written \emph{into} the state equation
rather than passed in as an external regressor whose relationship to $\Ttop$
the estimator is left to (mis)discover on its own:
\begin{equation}
  \mathrm d\Ttop = \tfrac{1}{C_{\mathrm{top}}}\Big[
      \bar P_{\mathrm{HP}}\,\delta(t) + \bar Q_{\mathrm{gshp}}\delta_g
      - \beta Q_{\mathrm{dist}} - \kappa(\Ttop-\Tbot) - UA_t(\Ttop-\Tamb)\Big]\mathrm dt
      + \sigma_t\,\mathrm d\omega_t,
  \label{eq:sde-top-switched}
\end{equation}
with $\delta(t)$ given by \eqref{eq:enable-law} (a state-dependent hysteresis
switch, structurally identical to how $\delta_g$ already enters via
$\bar Q_{\mathrm{gshp}}\delta_g$ in \S\ref{sec:gshp}) rather than the
free-floating $\alpha\Pel$ term of \eqref{eq:heat-inputs}. The only unknown
left in the ASHP term is $\bar P_{\mathrm{HP}}$, the effective heat rate while
enabled --- once the \emph{timing} ambiguity is removed by the known law, the
\emph{level} becomes a well-posed, identifiable parameter again, estimated
directly from how fast $\Ttop$ rises during each known-enabled stretch.

\paragraph{Caveats.}
\begin{enumerate}[nosep]
  \item This only identifies the dynamics \emph{within the operating band the
        thermostat has actually visited} (roughly \SIrange{31}{35}{\celsius},
        plus the transients into/out of it). It says nothing about the tank's
        response to conditions the old controller never explored --- which
        matters once the MPC starts commanding $\Tsp$ outside that band. A
        short designed identification test (step changes or a simple dithered
        sequence directly on $\Tsp$, independent of $\Ttop$) remains the clean
        long-term fix, and is required before trusting the model much beyond
        \SIrange{31}{35}{\celsius}.
  \item $\bar P_{\mathrm{HP}}$ is an \emph{average} effective rate over each
        enabled stretch. If the true in-band modulation depends on more than
        just enabled/disabled (e.g.\ how far below \SI{31}{\celsius} the
        stretch started, or ramp-up dynamics), that is extra structure not
        captured by \eqref{eq:enable-law} alone --- it only pins down
        \emph{when} the pump runs, not \emph{how hard}.
\end{enumerate}

\subsection{Stage 1 --- tank dynamics (both heat pumps together)}
Two states $x=[\Ttop,\Tbot]$, driven by the tank balances of \S\ref{sec:tank}
promoted to SDEs by adding a process-noise term to each:
\begin{align}
  \mathrm d\Ttop &= \tfrac{1}{C_{\mathrm{top}}}\!\Big[
      \alpha\Pel + \bar Q_{\mathrm{gshp}}\delta_g
      - \beta Q_{\mathrm{dist}} - \kappa(\Ttop-\Tbot) - UA_t(\Ttop-\Tamb)\Big]\mathrm dt
      + \sigma_t\,\mathrm d\omega_t, \label{eq:sde-top}\\
  \mathrm d\Tbot &= \tfrac{1}{C_{\mathrm{bot}}}\!\Big[
      \kappa(\Ttop-\Tbot) - UA_b(\Tbot-\Tamb)\Big]\mathrm dt
      + \sigma_b\,\mathrm d\omega_b. \label{eq:sde-bot}
\end{align}
The two tank sensors give the discrete observations
\begin{equation}
  y^{(1)}_k = \Ttop(t_k) + e^{(1)}_k, \qquad
  y^{(2)}_k = \Tbot(t_k) + e^{(2)}_k, \qquad
  e^{(j)}_k\sim\mathcal N(0,\varepsilon_j^2),
  \label{eq:tank-obs}
\end{equation}
mapped to columns \code{TT\_601} and \code{TT\_602}. The \emph{measured inputs}
$u=[\Pel,\ \delta_g,\ Q_{\mathrm{dist}},\ \Tamb]$ are supplied as known time series
($Q_{\mathrm{dist}}$ being the forecasted district load, not a flow); the
\emph{estimated parameters} are
\[
  \theta=\{\,C_{\mathrm{top}},\,C_{\mathrm{bot}},\,\kappa,\,UA_t,\,UA_b,\,
  \alpha,\,\beta,\,\bar Q_{\mathrm{gshp}},\,\sigma_t,\sigma_b,\,\varepsilon_1,\varepsilon_2\,\}.
\]
\texttt{ctsmTMB} runs the Kalman filter for a candidate $\theta$, reads off the
likelihood of the observed \code{TT\_601},\code{TT\_602} series, and maximises it
over $\theta$. Every parameter is pinned by a distinct feature of the data:
$\alpha$ by how fast $\Ttop$ rises per unit $\Pel$; $\beta$ by the $\Ttop$ droop per
unit forecast $Q_{\mathrm{dist}}$ during discharge (with $C_{\mathrm{top}}$ anchored
to the known volume so the two do not trade off); $UA_\bullet$ by the standing
cool-down when both pumps are off; $\kappa$ by how quickly $\Tbot$ follows $\Ttop$;
$\bar Q_{\mathrm{gshp}}$ as explained next.

\paragraph{GSHP: one constant, fitted here in Stage 1 --- there is no ``Stage 3''.}
The GSHP is \textbf{not} a separate stage and does \textbf{not} join the Stage-2
static equation \eqref{eq:staticmap}: that equation regresses power against a
\emph{modulating} set-point, and the on/off GSHP has none. Because it is pure on/off
with a nearly constant output, its whole model is a
\emph{single scalar} $\bar Q_{\mathrm{gshp}}$ (the heat it delivers while running),
switched by the measured binary $\delta_g=\code{PAC\_301\_RUN}\in\{0,1\}$ ---
the term $\bar Q_{\mathrm{gshp}}\delta_g$ in \eqref{eq:sde-top}. There is nothing to
curve-fit: flow$\times\Delta T$ or COP$\times$power can never be separated from
on/off data, and the tank cares only about the net heat, which is the one number
$\bar Q_{\mathrm{gshp}}$. It is identified \emph{for free} inside the same fit:
during the intervals when $\delta_g=1$, the tank heats faster than $\alpha\Pel$
alone can explain, and that unexplained extra rate of rise \emph{is}
$\bar Q_{\mathrm{gshp}}$. One number, pinned by one effect. All data are kept ---
when $\delta_g=0$ the term simply vanishes, so no GSHP-off filtering is needed.
(Only if the GSHP were ever metered would the $\alpha_g P_{\mathrm{el},g}$ form of
\eqref{eq:gshp} apply; for on/off it is unnecessary.)

\subsection{Stage 2 --- the actuator map (offline, all-measured)}
The nonlinear COP/$\Qmax$ map is needed \emph{only by the controller}, to predict
$\Pel$ (the cost) and the ceiling from a \emph{hypothetical} decision $\Tsp$ ---
which never occurs in the data. It is therefore fit \emph{separately and offline},
and it is easy because \textbf{every variable in it is measured}. Using
$\Ths\approx\Tsp$ at steady state and \eqref{eq:qhp-red}, the \emph{entire} static
actuator --- demand branch, capacity ceiling, and COP --- collapses into a
\textbf{single scalar equation}:
\begin{equation}
  \Pel = \frac{\,\overbrace{\min\!\big(\,G_h(\Tsp-\Thr),\ a_0+a_1\Tair\,\big)}^{\Qhp\ =\ \text{thermal demand, capped by ceiling}}\,}
             {\underbrace{b_0+b_1\Tair+b_2\Tsp}_{\mathrm{COP}(\Tair,\Tsp)}}
  \;\equiv\; g(x;\,p), \quad
  x=(\Tsp,\Thr,\Tair),\ \ p=(G_h,a_0,a_1,b_0,b_1,b_2).
  \label{eq:staticmap}
\end{equation}
The numerator is the delivered heat $\Qhp$: the thermal demand $G_h(\Tsp-\Thr)$
limited by the air-dependent ceiling $a_0+a_1\Tair$; dividing by the COP turns heat
into electrical power. All six parameters $p$ are fit from this one equation. The
GSHP is deliberately \emph{not} in it --- being on/off it carries no set-point to
regress against, so it degenerates to one constant and is folded into Stage~1
(above), not fitted here.

\paragraph{How a nonlinear static fit works.}
Each logged sample $i$ (a timestamp with the ASHP running) supplies measured
regressors $x_i=(\Tsp,\Thr,\Tair)_i$ and a measured output $y_i=P_{\mathrm{el},i}$. ``Static''
means the prediction $\hat y_i=g(x_i;p)$ is obtained by \emph{plugging the row's
numbers into \eqref{eq:staticmap}} --- no time integration, no dependence on any
other sample; each row is one independent point on a surface. Fitting chooses the
parameters $p$ that make the predictions sit closest to the measurements in the
least-squares sense:
\begin{equation}
  \min_{p}\ S(p) = \sum_{i=1}^{N}\big(y_i-g(x_i;p)\big)^2 .
  \label{eq:nls}
\end{equation}
This is the \emph{same} objective as fitting a straight line; the difference is only
\emph{how the minimum is found}. Because $p$ enters through a product/ratio, there
is no closed-form solution, so \eqref{eq:nls} is minimised \textbf{iteratively}:
start from a guess $p^{0}$, linearise $g$ about it through the sensitivities
$J_{ij}=\partial g(x_i;p)/\partial p_j$, solve a \emph{linear} least-squares problem
for a downhill correction $\Delta p$, update $p\leftarrow p+\Delta p$, and repeat
until $S$ stops decreasing (Gauss--Newton / Levenberg--Marquardt). ``Nonlinear least
squares'' is therefore just ordinary least squares reached by walking downhill; it
needs a sensible starting guess and, in principle, can settle in a local minimum ---
that is the only cost of the nonlinearity.

\paragraph{The capacity ceiling is fitted from saturation --- if the data reach it.}
There is no sensor for the ceiling $a_0+a_1\Tair$; it is identified from \emph{where
the machine stops responding}, and the $\min(\cdot)$ in \eqref{eq:staticmap} sorts
the samples into the two regimes that make this possible. \emph{Below} the ceiling
the compressor modulates: raising $\Tsp$ raises $\Pel$ and the supply reaches
set-point, so these rows ride the demand branch and pin $G_h$ and the COP $b_j$.
\emph{At} the ceiling the compressor is flat out: $\Pel$ plateaus and a
\emph{persistent} tracking error opens up ($\Ths<\Tsp$ although the machine is on).
Those saturated rows --- ``maxed out \emph{and} still short'' --- pin $a_0,a_1$, and
the fit places the ceiling at the \emph{kink} where the point cloud bends from
sloped to flat (equivalently, the upper envelope of $\Pel$ versus $\Tair$).

\emph{Caveat --- excitation.} The ceiling is identifiable \textbf{only if the data
actually hit it}. If the plant always kept headroom there are no saturated rows,
$a_0,a_1$ are not identifiable, and this is not a defect: the data hold no
information about a limit that was never reached. Then, for \emph{reproducing} the
data, drop the $\min$ and keep the demand branch alone; for the \emph{MPC} --- which
may command into saturation --- take the ceiling from the manufacturer's nameplate
capacity curve $\Qmax(\Tair)$ rather than the fit. Practical test: flag rows with
$\Pel$ pinned near its maximum \emph{and} $\Ths<\Tsp$; if enough of them span a range
of $\Tair$, fit $a_0,a_1$; otherwise fix them from the datasheet and identify only
$G_h$ and the COP.

\paragraph{Scale degeneracy --- why the two stages need each other.}
The Stage-2 regression has a scale ambiguity: multiplying $G_h$ and every $b_j$ by
a common factor leaves $\Pel$ unchanged. It therefore fixes only the \emph{shape}
of the COP's dependence on $(\Tair,\Tsp)$, not its absolute level. The absolute
scale is pinned by Stage~1's $\alpha$, because the tank thermal mass times the
measured temperature rise is an \emph{absolute} energy. Concretely:
\[
  \textbf{Stage 1 (tank)} \Rightarrow \text{absolute heat scale } \alpha;
  \qquad
  \textbf{Stage 2 (static regression)} \Rightarrow \text{conditions-dependence (shape)} .
\]
Anchor Stage~2 by fixing $b_0$, or by constraining its horizon-average COP to the
$\alpha$ obtained in Stage~1.

\subsection{Purist alternative --- one joint model}
The two stages can be fused into a single maximum-likelihood problem. Add the
logged set-point $\Tsp$ as a further input, and let the heat delivered be the
\emph{latent} (unmeasured) quantity computed from the decision variable, capped by
the ceiling:
\begin{equation}
  \Qhp = \min\!\big(\,G_h(\Tsp-\Tbot),\ a_0+a_1\Tair\,\big).
  \label{eq:qhp-latent}
\end{equation}
This same $\Qhp$ is used in \emph{two} places at once. First it drives the tank top
node, i.e.\ $\alpha\Pel$ in \eqref{eq:sde-top} is replaced by $\Qhp$:
\begin{equation}
  \mathrm d\Ttop = \tfrac{1}{C_{\mathrm{top}}}\!\Big[
      \Qhp + \bar Q_{\mathrm{gshp}}\delta_g
      - \beta Q_{\mathrm{dist}} - \kappa(\Ttop-\Tbot) - UA_t(\Ttop-\Tamb)\Big]\mathrm dt
      + \sigma_t\,\mathrm d\omega_t .
  \label{eq:sde-top-joint}
\end{equation}
Second, it appears in a \emph{third observation equation} against the measured
electrical power (\code{PAC\_501\_PUISSANCE}):
\begin{equation}
  y^{(3)}_k = \frac{\Qhp(t_k)}{\,b_0+b_1\Tair+b_2\Tsp\,} + e^{(3)}_k,
  \qquad e^{(3)}_k\sim\mathcal N(0,\varepsilon_3^2).
  \label{eq:pel-obs}
\end{equation}
Because the tank observation \eqref{eq:tank-obs} and the power observation
\eqref{eq:pel-obs} now share the \emph{same} latent $\Qhp$, the scale degeneracy of
Stage~2 is resolved automatically: the tank fixes the absolute size of $\Qhp$ and
the power equation fixes how it splits into $G_h$ and the COP denominator. One MLE
then returns $\{C_\bullet,\kappa,UA_\bullet,\beta,\bar Q_{\mathrm{gshp}},G_h,a_0,a_1,
b_0,b_1,b_2,\sigma_\bullet,\varepsilon_\bullet\}$ --- tank \emph{and} actuator
parameters, mutually consistent.

\textbf{Trade-off.} This is elegant but harder to initialise and debug: more local
optima, the non-smooth $\min(\cdot)$ kink in \eqref{eq:qhp-latent}, and $\alpha$ is
now \emph{implied} by $G_h$ and the COP rather than fit as its own number.
\textbf{Recommendation:} run the two-stage version first, then --- only if the two
stages disagree --- move to the joint model, using the two-stage estimates as its
starting guess.

\subsection{Remaining steps}
\begin{enumerate}[nosep]
  \item \textbf{Thermocline (optional):} calibrate the sensor-crossing heights and
        $g(\theta)$ for the $\Vtop$ observer \eqref{eq:vtop}; trigger it by the
        opposite-curvature diagnostic of \S\ref{sec:tank}.
  \item \textbf{Validation:} multi-step-ahead ($k$-step) prediction on held-out
        weeks (not one-step filtering, which hides model error), inspecting the
        $\Ttop$ fit specifically around the discharge ``cliffs''.
\end{enumerate}

\paragraph{Data hygiene.}
Time index is UTC; cumulative counters must be differenced to rates; report and
handle the known 15-minute sampling gaps (\texttt{ctsmTMB} accepts non-uniform
timestamps, but the gaps must be genuine, not silently interpolated); keep the
inspection cache faithful (no resampling before identification).

\end{document}
